\subsection{一致空间的逆极限}

设 I 是一个偏序集，其偏序记作 \(\alpha \leqslant \beta\)。对每个 \(\alpha \in I\)，设 \(X_{\alpha}\) 是一致空间；又对每对指标 \(\alpha, \beta\)（满足 \(\alpha \leqslant \beta\)），设 \(f_{\alpha \beta}\) 是 \(X_{\beta}\) 到 \(X_{\alpha}\) 的映射。

我们称 \((\mathbf{X}_{\alpha}, f_{\alpha\beta})\) 是一致空间的逆系统，如果 (i) \((\mathbf{X}_{\alpha}, f_{\alpha\beta})\) 是集合的逆系统（参见第一章附录第 1 目），且 (ii) 只要 \(\alpha \leqslant \beta\)，\(f_{\alpha\beta}\) 就是一致连续的。在 \(X = \varprojlim X_{\alpha}\) 上，使诸典范映射 \(f_{\alpha}: X \to X_{\alpha}\) 都一致连续的最粗一致结构，称为（关于 \(f_{\alpha\beta}\) 的）诸 \(X_{\alpha}\) 的一致结构的逆极限；赋有这一最粗一致结构的集合 X 称为一致空间的逆系统 \((\mathbf{X}_{\alpha}, f_{\alpha\beta})\) 的逆极限。第一章 § 4 第 4 目中建立的拓扑空间逆极限的所有性质（命题 9 除外）在把“拓扑”换成“一致结构”、“连续映射”换成“一致连续映射”之后仍然成立。此外：

命题 10 设 I 是有向集，\((\mathbf{X}_{\alpha}, f_{\alpha \beta})\) 是以 I 为指标的一致空间逆系统，J 是 I 的共尾子集。对每个 \(\alpha \in I\)，设 \(f_{\alpha}\) 是 \(\mathbf{X} = \varprojlim \mathbf{X}_{\alpha}\) 到 \(\mathbf{X}_{\alpha}\) 的典范映射，并用 \(g_{\alpha}\) 表示 \(f_{\alpha} \times f_{\alpha}\)。则集合族 \(\overline{g}_{\alpha}^{1}(\mathrm{V}_{\alpha})\)（其中 \(\alpha\) 取遍 J，且对每个 \(\alpha \in J\)，\(\mathrm{V}_{\alpha}\) 取遍 \(\mathrm{X}_{\alpha}\) 的一个一致邻域基本系）构成 X 的一致邻域基本系。

证明留给读者；它只需把第一章 § 4 第 4 目命题 9 的证明直接加以改编。

最后，在 \(X = \varprojlim X_{\alpha}\) 上由诸 \(X_{\alpha}\) 的一致结构的逆极限所诱导的拓扑，就是诸 \(X_{\alpha}\) 的拓扑的逆极限。

